\documentclass[11pt]{article}
\usepackage[a4paper,margin=24mm]{geometry}
\usepackage[T1]{fontenc}
\usepackage{lmodern,amsmath,amssymb,amsthm,mathtools,mathrsfs}
\usepackage[hypertexnames=false,hidelinks]{hyperref}
\setlength{\emergencystretch}{2em}
\newcommand{\C}{\mathbb C}\newcommand{\R}{\mathbb R}
\newcommand{\PP}{\mathcal P}\newcommand{\dd}{\,d}
\newcommand{\norm}[1]{\left\lVert#1\right\rVert}
\newtheorem{theorem}{Theorem}\newtheorem{lemma}[theorem]{Lemma}
\title{Full-window source products and the retained class:\\finite quotient comparisons in the original arithmetic metric}
\author{Split-Zero research programme}
\date{19 September 2026}
\begin{document}\maketitle

The user-supplied continuation \cite{incoming} proposes an evaluation of the
two complete source-response products and the four-cutoff terminal-class
return. We prove these evaluations below, retaining the covariance correction
as an exact finite sum. In particular, their proof does not require a new
uniform exponential estimate for that covariance. The already established
finite original-source multiplication and exterior estimates suffice, at
every fixed primary multiplicity. We also give the exact receiving identity
for the same class in the original observation quotient. Statements about
that observation keep its actual projection angle and complex responses.

\section{Original packet, attained class and source energies}
Fix the actual quartet data of the preceding calculation:
\begin{equation}\tag{FW1}
 \begin{gathered}
 \rho=\tfrac12+\delta+i\gamma,\quad 0<\delta<\tfrac12,\quad
 \gamma>2,\quad m_0\ge1,\\
 k\equiv1\pmod4,\quad k\ge5,\quad
 e=1+k(m_0-1),\quad q=e(k+1)^2.
 \end{gathered}
\end{equation}
Keep the full quartet factor $h$ in $2\xi$, its full arithmetic unit and
the literal measure and physical coordinate
\[
 w_h(y)=|(2\xi/h)(1/2+iy)|^2/(2\pi),\quad
 dm=w_h^{*k}(y)\dd y,\quad S=c+iy,\quad c=k/2.
\]
The original complete polynomial and terminal class are
\begin{equation}\tag{FW2}
 \begin{gathered}
 \chi(S)=\prod_{a,b=0}^k[S-c-(2a-k)\delta-i(2b-k)\gamma]^e,
 \quad E=\C[S]/(\chi),\\
 \lambda=k\rho=c+d+it,\quad d=k\delta,\quad t=k\gamma,
 \quad z=t-id,\quad R=|z|=k\sqrt{\delta^2+\gamma^2},\\
 \chi_\lambda=\chi/(S-\lambda)^e,\quad b_\lambda=1/\chi_\lambda(\lambda),
 \qquad v=[b_\lambda\chi/(S-\lambda)].
 \end{gathered}
\end{equation}
No root difference, repeated terminal power or arithmetic scalar has been
changed. Write $Q(y)=i^{-q}\chi(c+iy)$. Conjugation and negation preserve
its full root multiset, and no root is real. Thus $Q$ is real, even,
strictly positive on the real line, and $Q(z)=Q(\bar z)=0$.
Put $dW=Q^2dm$. Its real monic orthogonal polynomials and squared norms
are $U_j,\nu_j$. Define
\begin{equation}\tag{FW3}
 C_j(z)=\int\frac{U_j(y)}{z-y}\dd W(y),\qquad
 H_r(z)=\min_{P\in\C[y]_{<r}}\int\left|\frac1{z-y}-P(y)\right|^2\dd W(y).
\end{equation}
The minimum is attained by $\sum_{j<r}C_j(z)U_j/\nu_j$.
Multiplication by $-b_\lambda i^{q-1}Q$ maps this affine fibre exactly
to all representatives of $v$ through degree $N=q-1+r$; its inverse
divides by that same specified nonzero factor. This gives
\begin{equation}\tag{FW4}
 E_N=v^*G_Nv=|b_\lambda|^2H_r,
 \quad G_N=(J_NH_N^{\rm source,-1}J_N^*)^{-1},
 \quad H_r-H_{r+1}=|C_r|^2/\nu_r.
\end{equation}
Here $J_N$ is the original remainder map and $H_N^{\rm source}$ is the
Gram of $1,S,\ldots,S^N$ in $dm$. Every original relation is included
in its minimum. The finite differences in (FW4) need no infinite
orthogonal expansion.

Let $p_n(S)$ be the original $S$-monic source polynomial,
$\omega_n=\norm{p_n}_m^2$, and $b_n=[p_n]$. In the three columns
$(b_N,b_{N+1},v)$ put $F_{ij}=W_i^*G_NW_j$ and
\[
 p_{1,N}=\frac{|F_{13}|^2}{F_{11}E_N},\qquad
 p_{2,N}=\frac{|F_{23}|^2}{F_{22}E_N},\qquad
 \Delta_r=\omega_{q+r}/\nu_r.
\]
The full monic relation of degree $q+r$ is $\chi\pi_r$, where
$\pi_r(c+iy)=i^rU_r(y)$. Its difference from $p_{q+r}$ has degree
at most $q+r-1$, represents $b_{q+r}$, and is orthogonal to every
preceding relation. Squaring its norm gives $F_{22}=\nu_r-\omega_{q+r}$.
The same monic projection in degree $q+r-1$ gives
\begin{equation}\tag{FW5}
 p_{2,q-1+r}=\frac{H_r-H_{r+1}}{(1-\Delta_r)H_r},\qquad
 p_{1,q-1+r}=\frac{\Delta_{r-1}}{1-\Delta_{r-1}}
                  \left(\frac{H_{r-1}}{H_r}-1\right)\quad(r\ge1).
\end{equation}
At $r=0$ the unchanged first formula is $p_{1,q-1}=\omega_{q-1}/H_0$.
These are precisely the full minimum projections in RBE12--19
\cite{rbe}; the changes of variable above explain every factor of $i$.

\begin{lemma}[Finite residual Wronskian]
For $r\ge1$,
\begin{equation}\tag{FW6}
 H_r=\frac{\operatorname{Im}(C_r\overline{C_{r-1}})}{d\nu_{r-1}}>0.
\end{equation}
\end{lemma}
\begin{proof}
Evenness and monicity give $U_{j+1}=yU_j-(\nu_j/\nu_{j-1})U_{j-1}$.
Integrating the recurrence gives $C_1=zC_0-\nu_0$ and
$C_{j+1}=zC_j-(\nu_j/\nu_{j-1})C_{j-1}$ for $j\ge1$.
Since $\operatorname{Im}C_0=dH_0$, the right side for $r=1$ is
$H_0-|C_0|^2/\nu_0=H_1$. At each later step its difference is
$-|C_j|^2/\nu_j$, exactly (FW4). Positivity follows because the
rational function cannot equal a polynomial on an interval.
\end{proof}
The recurrence, simple real zeros and interlacing used below are the
classical orthogonal-polynomial results recorded in DLMF
\cite[\S18.2(iv), (vi), (vii)]{DLMF}. The residual identities themselves
have just been proved for the original arithmetic measure.

\section{Finite control of the central interval and rational minima}
The original CAI16--25 estimates \cite{cai} state, for every finite degree
$n$, with $M=21+4m_0$,
\begin{equation}\tag{FW7}
 \frac{a_k}{[1+4(n+M)^2]^M}\norm P_\sigma^2
 \le\norm P_m^2\le A_k\norm P_\sigma^2,
 \qquad \norm{yP}_m\le C_{\rm mult}(h)(2n+\lfloor k/3\rfloor)\norm P_m.
\end{equation}
Here $\sigma=|\Gamma(1/4+iy/2)|^2/(2\pi)$, of mass $\sqrt{2\pi}$;
the positive constants $a_k,A_k,C_{\rm mult}$ retain the source masses
and quartet from CAI16 and CAI22. In particular
$\log(A_k/a_k)=O_h(k+\log k)$. These estimates hold for complex
polynomials, not merely real monomials. We use $n=2q+2$ and denote the
two comparison constants by $\ell,u$, and $\Lambda=\log(u/\ell)$.
For Gamma, $\norm{yP}_\sigma\le2(n+1)\norm P_\sigma$, from its exact
three-term recurrence. No arithmetic metric is replaced by this reference.

Choose a fixed $\epsilon_*\le2^{-16}$ with $\epsilon_*^2$ below the
support of $\rho_{2,\pi}$, put $Y=\epsilon_*q$, $D=q+2$, and impose
$R\le Y/4$. With the original Gamma lower-envelope constant $c_G>0$,
set
\begin{equation}\tag{FW8}
 \eta=\frac{u}{\ell}\frac{\sqrt{2\pi}}{c_G}
 \frac{(D+1)^2\sqrt{1+2q}}q,100^D e^{\pi q}
 \left(\frac{Y^2+R^2}{q^2-R^2}\right)^q.
\end{equation}
Then, for every $\deg P\le D$,
\begin{equation}\tag{FW9}
 \int_{|y|\le Y}|P|^2\dd W\le\eta\norm P_W^2,
 \qquad \eta\le e^{-q\log16+O_h(k+\log q)}.
\end{equation}
To verify this, the shifted Legendre basis on $[q,2q]$ gives
\[
 \sup_{|y|\le Y}|P(y)|^2\le (D+1)^2 100^Dq^{-1}\int_q^{2q}|P|^2;
\]
its argument belongs to $[-4,4]$, and the defining Legendre formula
bounds each degree-$j$ basis value by $10^j$ times its normalization.
Opposite-root pairs give $|Q|^2\le(Y^2+R^2)^q$ centrally and
$|Q|^2\ge(q^2-R^2)^q$ on $[q,2q]$. The central measure has mass at
most $u\sqrt{2\pi}$, by (FW7) for $1$; the full denominator is at
least $\ell\int_q^{2q}|PQ|^2\sigma$. The Gamma lower density there
is $c_G(1+2q)^{-1/2}e^{-\pi q}$. These statements prove (FW8) without
assuming a pointwise comparison of arithmetic and Gamma densities.
Finally $100e^\pi(17\epsilon_*^2/16)/(1-\epsilon_*^2/16)<1/16$,
which proves the last assertion. Both complementary regions remain in
the original forms.

Put $W^\circ=1_{|y|>Y}W$, and use a superscript $\circ$ for its
polynomials, norms and residuals. Take
$C=C_{\rm mult}(h)[2(2q+1)+\lfloor k/3\rfloor]$ and
$\xi=\eta(C+R)^2/d^2$. For $0\le r\le q+2$,
\begin{equation}\tag{FW10}
 (1-\xi)H_r\le H_r^\circ\le H_r,\qquad
 |T_r-T_r^\circ|\le\frac{\xi}{1-\xi},\quad T_r=1-H_{r+1}/H_r.
\end{equation}
The second statement uses $r\le q+1$. Indeed each candidate has
numerator $Z=1-(z-y)P$ of degree at most $r$. On the central interval
its squared norm is at most $d^{-2}\int|Z|^2\dd W$ restricted there.
Apply (FW9) to $Z$. The polynomial $F=QZ/(z-y)$ exists because $Q(z)=0$,
and has degree at most $q+r-1$. Moreover
$\norm Z_W=\norm{(z-y)F}_m\le(C+R)\norm F_m$.
This proves the candidate-wise lower bound and hence the attained one.
Writing $H_r^\circ=a_rH_r$, $1-\xi\le a_r\le1$, proves the ratio
bound in (FW10). Degree $2q+2$ in (FW7) includes its final output.

For $j\le q+1$ define
\[
 M_j^\circ=\int\frac{(U_j^\circ)^2}{y^2}\dd W^\circ,
 \qquad \widehat C=\frac{C_{\rm mult}(h)[2(2q+1)+\lfloor k/3\rfloor]}{\sqrt{1-\eta}}.
\]
Multiplication on $QU_j$ and (FW9) bound
$\norm{yU_j^\circ}_{W^\circ}\le\widehat C\norm{U_j^\circ}_{W^\circ}$.
Cauchy--Schwarz applied to $U_j/y$ and $yU_j$, and the pointwise
$|y|>Y$, now give
\begin{equation}\tag{FW11}
 \widehat C^{-2}\le M_j^\circ/\nu_j^\circ\le Y^{-2},\qquad
 \nu_j^\circ/\nu_{j-1}^\circ\le\widehat C^2.
\end{equation}
The second inequality is the minimum monic norm applied to $yU_{j-1}$.

\section{Consecutive amplitudes and exact product cancellation}
Write $U_{2j}^\circ(y)=P_j(y^2)$, $U_{2j+1}^\circ(y)=yT_j(y^2)$.
Orthogonality yields
\begin{equation}\tag{FW12}
 \int T_j(y^2)\dd W^\circ=\frac{\nu_{2j}^\circ}{P_j(0)},\quad
 M_{2j+1}^\circ=\frac{T_j(0)\nu_{2j}^\circ}{P_j(0)},\quad
 P_{j+1}(0)=-\frac{\nu_{2j+1}^\circ}{\nu_{2j}^\circ}P_j(0).
\end{equation}
For the first identity, $P_j-P_j(0)$ is $x$ times a degree-$j-1$
polynomial, so its product with $T_j$ integrates to zero; while
$\int P_jT_j=\int P_j^2=\nu_{2j}^\circ$. For the second,
$T_j-T_j(0)$ is likewise $x$ times a lower polynomial, orthogonal to
$T_j$ in the $x$-weighted measure. The last equality is the monic
recurrence at zero. Quadratic transformation places every positive
squared zero above $Y^2$, and the two consecutive zero lists interlace.

Put $\varrho=R^2/Y^2$. Orthogonality of
$(U_j(z)-U_j(y))/(z-y)$ against $U_j$ proves the exact identity
\begin{equation}\tag{FW13}
 U_j^\circ(z)C_j^\circ(z)
 =\int\frac{(U_j^\circ(y))^2}{z-y}\dd W^\circ
 =-zM_j^\circ A_j(z),\qquad |A_j(z)-1|\le\frac{\varrho}{1-\varrho}.
\end{equation}
The last step pairs $y,-y$: $A_j$ is the probability-weighted average
of $(1-z^2/y^2)^{-1}$ with weight $(U_j^\circ)^2/y^2$.
For interlacing positive squared zeros, the signed counting function
has absolute value at most one. Integrating
$\partial_x\log(1-z^2/x)$ from $Y^2$ to infinity therefore bounds
the logarithm of the consecutive root-product ratio by
$\int_{Y^2}^\infty R^2/[x(x-R^2)]\dd x=-\log(1-\varrho)$.
This proves the complex ratio estimate, including its phase and
without a degree factor. In conjunction with (FW12)--(FW13) it gives
\begin{equation}\tag{FW14}
 \frac{C_r^\circ}{C_{r-1}^\circ}=
 \begin{cases}
 -z(M_r^\circ/\nu_{r-1}^\circ)(1+v_r),&r\text{ even},\\
 z^{-1}(\nu_{r-1}^\circ/M_{r-1}^\circ)(1+v_r),&r\text{ odd},
 \end{cases}
 \qquad |v_r|\le8\varrho\quad(\varrho\le1/16).
\end{equation}
For explicit verification of the constant, each $A$-ratio differs from
one by at most $2\varrho/(1-2\varrho)$, and the product ratio by at
most $\varrho/(1-\varrho)$; multiplying the factors gives a bound
smaller than $8\varrho$ on the stated interval.

Define positive finite quantities
\begin{equation}\tag{FW15}
 b_r=\begin{cases}
 R^2M_r^\circ/\nu_r^\circ,&r\text{ even},\\
 (\nu_{r-1}^\circ)^2/(\nu_r^\circ M_{r-1}^\circ),&r\text{ odd},
 \end{cases}\quad
 a_*=32(1+R/d)\varrho+
       \frac{\xi}{1-\xi}\frac{\widehat C^2}{R^2}.
\end{equation}
Impose the explicit eventual guards $\eta,\xi<1/4$, $\varrho\le1/16$,
$32(1+R/d)\varrho\le1/4$, and $a_*\le1/4$. Then
\begin{equation}\tag{FW16}
 T_r=b_r(1+\varepsilon_r),\quad |\varepsilon_r|\le a_*,
 \qquad p_{2,q-1+r}=\frac{b_r(1+\varepsilon_r)}{1-\Delta_r},
 \quad 0\le r\le q+1.
\end{equation}
For $r\ge1$ substitute (FW14) in (FW6). The expression is
\[
 T_r^\circ=\frac{d\nu_{r-1}^\circ|C_r^\circ/C_{r-1}^\circ|^2}
 {\nu_r^\circ\operatorname{Im}(C_r^\circ/C_{r-1}^\circ)}.
\]
The relative error of the imaginary part is at most $8(R/d)\varrho$;
the numerator error is at most $16\varrho+64\varrho^2$.
Their ratio is within $32(1+R/d)\varrho$ of one on the guard.
For $r=0$, use $C_0=-zM_0 A_0$, $H_0=\operatorname{Im}C_0/d$,
which has the same bound. Finally (FW11) gives
$b_r\ge R^2/\widehat C^2$ for even $r$ and
$b_r\ge Y^2/\widehat C^2\ge R^2/\widehat C^2$ for odd $r$.
Thus (FW10) gives exactly the remaining relative allowance in $a_*$.

Let $w_0=w_q=1$, $w_r=2$ for $0<r<q$, so $\sum w_r=2q$.
Define the exact positive covariance contribution
\begin{equation}\tag{FW17}
 D_k=\sum_{r=0}^q w_r[-\log(1-\Delta_r)],\qquad
 \mathscr A_i=-\sum_{r=0}^q w_r\log p_{i,q-1+r}.
\end{equation}
Equation (FW12) and the two definitions of $b$ imply
$b_{2j}b_{2j+1}=R^2\nu_{2j}^\circ/\nu_{2j+1}^\circ$ and
$|U_q^\circ(0)|=\prod_{j<q/2}\nu_{2j+1}^\circ/\nu_{2j}^\circ$.
Taking the literal endpoint-weighted sum therefore proves
\begin{equation}\tag{FW18}
 \boxed{\mathscr A_2=
 2\log|U_q^\circ(0)|-2q\log R
 -\log\frac{M_q^\circ/\nu_q^\circ}{M_0^\circ/\nu_0^\circ}
 -D_k+\mathcal E_k,\qquad |\mathcal E_k|\le4qa_*.}
\end{equation}
Every inverse moment except the displayed endpoint ratio has cancelled.
Its logarithm is bounded by $2\log(\widehat C/Y)$.

\paragraph{A finite covariance bound suffices.}
Let $L_{\rm PT}=\delta q(k+1)/2$ and
$B=C_{\rm mult}(h)^2(4q+\lfloor k/3\rfloor)^2$.
The original exterior determinant bound gives $\kappa_N\ge L_{\rm PT}$
at every $N\ge q-1$. Direct substitution of the two monic boundary
norms gives
$\kappa_{q-1}^2=(\omega_q/\omega_{q-1})(\Delta_0^{-1}-1)$ and
$\kappa_N^2=(\omega_{N+1}/\omega_N)(1-\Delta_{r-1})(\Delta_r^{-1}-1)$
at $r\ge1$. The source multiplication bound controls the recurrence
coefficient by $B$. Consequently
\begin{equation}\tag{FW19}
 0<\Delta_r\le\frac{B}{B+L_{\rm PT}^2},\qquad
 0\le D_k\le2q\log(1+B/L_{\rm PT}^2)=O_h(q/k^2)=O_h(k).
\end{equation}
For $r=q+1$, enlarge $4q$ in $B$ to $4q+2$; the same proof applies.
This derives exactly the needed estimate from PT7, PT23 and CAI25,
as in RBE26. There is no use of a conjectured covariance limit.
Moreover $qa_*=O_h(k^2/q)+e^{-cq+O_h(k+\log q)}$ and
$\widehat C/Y=O_h(1)$. Thus the exact correction in (FW18) fits in
$O_h(k+\log q)$ at every fixed $m_0$.

\section{Two quotient metrics evaluate the polynomial value}
Push $W^\circ$ forward under the literal $x=(y/q)^2$, preserving its
full mass. Put $n=q/2$, and write
$U_q^\circ(y)=q^q\widehat P_n((y/q)^2)$. For both this arithmetic
measure and the Gamma reference with the same full $Q$ and gap, put
\[
 D_n^\mu(a)=\det[\int x^{i+j+a}\dd\mu]_{i,j<n},\qquad
 f_\mu(a)=\log D_n^\mu(a).
\]
Solving the monic orthogonality equations for the constant coefficient
by Cramer's rule yields $|\widehat P_n(0)|=D_n^m(1)/D_n^m(0)$.
All determinants are positive. Their squared-Vandermonde integral,
with its factor $1/n!$, makes $f_\mu''(a)$ the variance of
$\sum\log x_j$. Differentiation is justified by the positive lower
endpoint and exponential tail. Hence
\begin{equation}\tag{FW20}
 \frac{f_m(0)-f_m(-2)}2\le\log|\widehat P_n(0)|
 \le\frac{f_m(2)-f_m(0)}2.
\end{equation}

Here is the full metric comparison. On $V=\C[x]_{<n}$ and
$xV$, their common subspace is $x\C[x]_{<n-1}$. In each determinant
factor it out by a Schur complement. The common Gram is identical
within each measure and cancels. The remaining factors are the
attained norms of $1$ and $x^n$ modulo this same subspace.
An inequality $\ell^\circ\norm F_\Gamma^2\le\norm F_m^2
\le u^\circ\norm F_\Gamma^2$ passes to each affine minimum.
The ratio of the two quotient factors is therefore between
$\ell^\circ/u^\circ$ and $u^\circ/\ell^\circ$ times its Gamma value.
No determinant dimension appears. Applied to (FW7)--(FW9), this gives
\begin{equation}\tag{FW21}
 |[f_m(2)-f_m(0)]-[f_\Gamma(2)-f_\Gamma(0)]|
 \le\Lambda^\circ,\quad
 \ell^\circ=(1-\eta)\ell,\quad
 u^\circ=u/(1-\eta_\Gamma),\quad
 \Lambda^\circ=\log(u^\circ/\ell^\circ).
\end{equation}
The Gamma central fraction $\eta_\Gamma$ is (FW8) with $u/\ell=1$.
Every polynomial here, after multiplication by $Q$ and the specified
substitution, has degree at most $2q$.

For the negative slope use $V,x^{-1}V$, with common subspace
$\C[x]_{<n-1}$. To compare their entire sum, write each rational
element as $F/x$ with $\deg F\le n$. In either gapped measure,
the full polynomial multiplication bound, including $Q$, yields
$\norm{yF}_\mu\le B_\mu\norm F_\mu$, with
$B_m,B_\Gamma=O_h(q)$ explicitly given by their degree-$2q$ constants
divided by $\sqrt{1-\eta_\mu}$.
Jensen's inequality for $t\mapsto t^{-2}$ in the probability measure
$|F|^2\dd\mu/\norm F_\mu^2$ gives
\[
 (q/B_\mu)^4\norm F_\mu^2\le\norm{F/x}_\mu^2
 \le\epsilon_*^{-4}\norm F_\mu^2.
\]
Combining these two estimates with (FW21)'s polynomial comparison
gives a rational logarithmic width
\[
 \Lambda^{\rm rat}=\Lambda^\circ+
     4\log\frac{B_mB_\Gamma}{\epsilon_*^2q^2}=O_h(k+\log q).
\]
The identical common-Gram cancellation proves
\begin{equation}\tag{FW22}
 |[f_m(0)-f_m(-2)]-[f_\Gamma(0)-f_\Gamma(-2)]|
 \le\Lambda^{\rm rat}.
\end{equation}
In particular the following finite interval has completely specified
Gamma endpoints:
\begin{equation}\tag{FW23}
 \frac{f_\Gamma(0)-f_\Gamma(-2)-\Lambda^{\rm rat}}2
 \le\log|\widehat P_n(0)|\le
 \frac{f_\Gamma(2)-f_\Gamma(0)+\Lambda^\circ}2.
\end{equation}

\section{Gamma slopes, both source products and the arithmetic correction}
Write $L_{\rm eq}=\int\log x\dd\rho_{2,\pi}$ and
$F=-\min\mathcal E_{2,\pi}$, with the original equilibrium conventions.
The complete Gamma calculation AG10--18 \cite{ag} retains both
complementary integration regions, all roots and the fixed potential
$V(x)=\pi\sqrt x-2\log x$. At $n=q/2$ the exact density factors for
$a=-2,0,2$ are
\begin{equation}\tag{FW24}
 e^{-nV(x)+W_{a,k}(x)},\quad
 W_{a,k}=(a-3/4)\log x-\tau_k/x,\qquad
 \tau_k=\frac{(\gamma^2-\delta^2)k(k+2)}{3q}.
\end{equation}
On the fixed positive interval the complete root remainder has
supremum $O_h(k^4/q^3)$, so its determinant cost is $O_h(k^4/q^2)=O_h(1)$.
The Gamma-density remainder costs $O(q^{-1})$. Outside that interval
the complete-form bounds of AG11--14 are exponentially small.
For the three fixed powers here the same bounds apply: on the retained
gap $x^{-2}\le\epsilon_*^{-4}$, while at infinity $x^2$ adds only two
fixed degrees to the polynomial estimate. The two additional degrees
were included in (FW7). Constants and small-cutoff guards can thus be
enlarged before, not after, the determinant comparison.

Borot and Guionnet's one-cut analytic-statistic theorem
\cite[Propositions 1.2 and 5.2]{BG}, in the monic determinant convention
proved in AG18, gives
\[
 \log D_n[e^{-nV+W}]-\log D_n[e^{-nV}]
       =n\int W\dd\rho_{2,\pi}+O_W(1).
\]
Its hypotheses hold on that fixed positive interval: AG6 gives a positive
one-cut density, nonzero edge factors and a strict exterior effective-potential
gap. Analyticity at zero is not invoked. Uniformity in the actual bounded
$\tau_k$ follows by convexity: subtract the displayed linear statistic
from the logarithm and bound this convex function at $-2T,0,2T$, with
$T>1+\sup|\tau_k|$. Convexity bounds it above on $[-T,T]$ and below by
expressing zero as a convex combination of an interior point and the
opposite endpoint. This gives a common finite bound, separately for each
of the three fixed powers. We have therefore proved
\begin{equation}\tag{FW25}
 f_\Gamma(2)-f_\Gamma(0)=qL_{\rm eq}+O_h(\log q),\qquad
 f_\Gamma(0)-f_\Gamma(-2)=qL_{\rm eq}+O_h(\log q).
\end{equation}
Equations (FW23)--(FW25), including the exact scaling of the monic
polynomial, give
\begin{equation}\tag{FW26}
 \log|U_q^\circ(0)|=q\log q+(q/2)L_{\rm eq}+O_h(k+\log q).
\end{equation}

\begin{theorem}[Complete original source-response products]
At every fixed original multiplicity, as $k\to\infty$ along (FW1),
\begin{align}\tag{FW27}
 \mathscr A_2&=2q\log\frac{q}{k\sqrt{\delta^2+\gamma^2}}
                    +qL_{\rm eq}+O_h(k+\log q),\\
 \mathscr A_1&=J_k^{\rm action}-2q\log(qk\sqrt{\delta^2+\gamma^2})
                    -qL_{\rm eq}+O_h(k+\log q).\tag{FW28}
\end{align}
\end{theorem}
\begin{proof}
(FW18), (FW19) and (FW26) prove the first equality. For the second,
the original full-action boundary identity RBE3--7 and EC give
$p_{1,N}p_{2,N}=|\mathfrak c_N|^2/\kappa_N^2$ and
$\mathfrak c_N=d+i(-1)^r(t-\varepsilon_N)$ with
$0<\varepsilon_N=O_h(k^3/q^2)$ throughout the window.
Thus $\log|\mathfrak c_N|=\log R+O_h(k^2/q^2)$, retaining its signed
phase before taking its modulus. The original action is
$J_k^{\rm action}=2\sum w_r\log\kappa_{q-1+r}$. Hence exactly
\[
 \mathscr A_1+\mathscr A_2=J_k^{\rm action}-4q\log R
                    +O_h(k^2/q).
\]
Subtract (FW27). All errors are $o_h(q/\log q)$ because
$q=[1+k(m_0-1)](k+1)^2$ with fixed $m_0$.
\end{proof}

In the unchanged constants of AG26--29, put
$C_\Gamma=2L_{\rm eq}-8\log2+4$,
$C_B=9-8\log2+F$, $C_\sigma=F+6-3L_{\rm eq}/2-6\log2+2\log\pi$,
$A_J=C_\sigma+8\log2-4-(4m_0-2)C_\Gamma$ and
$r_h=2(\gamma^2-\delta^2)M_{-1}/3$, where
$M_{-1}=\int x^{-1}\dd\rho_{2,\pi}$.
Retain $I_h=\int_0^{\pi/2}[(\log M_h)'(\theta)]^2\dd\theta$ with
$M_h(\theta)=\int e^{\theta y}w_h(y)\dd y$, and
$c_\zeta=2\gamma_{\rm E}-\log(2\pi)$.
Substitution of the proved action AG29, with its SR1--48 finite remainder,
in (FW28) yields
\begin{equation}\tag{FW29}\begin{aligned}
 \mathscr A_1={}&C_Bq^2+2q\log(q/k)
 +(A_J-L_{\rm eq}-2\log\sqrt{\delta^2+\gamma^2})q\\
 &+\frac{\log2}{\pi}I_hk^2-r_hk(k+2)
 +\frac{C_\Gamma q}{2(\log q+c_\zeta)}+o_h(q/\log q).
\end{aligned}\end{equation}
This use of the action is a receiving substitution, not an independent
reproof of its arithmetic asymptotic. In the Gamma source the same
proof of (FW27) gives the identical leading product; consequently
\begin{equation}\tag{FW30}
 \mathscr A_2^m-\mathscr A_2^\Gamma=O_h(k+\log q),\qquad
 \mathscr A_1^m-\mathscr A_1^\Gamma
 =\delta_{k,1}^{\rm ar}+O_h(k+\log q).
\end{equation}
Here $\delta_{k,1}^{\rm ar}$ is the original arithmetic baseline
difference; the action's monic-window and penalty comparisons contribute
only $O_h(k+\log q)$, as in CAI and AG28--29. Thus the $q/\log q$
arithmetic term is assigned to a specified source-correlation product.

\section{The retained class's actual four-cutoff norm}
The auxiliary zero argument is used only in $W^\circ$, whose support
avoids zero. For even $r=2j$, the polynomial
$(U_r^\circ(y)/U_r^\circ(0)-1)/y$ has degree less than $r$ and leaves
residual $-U_r^\circ(y)/(yU_r^\circ(0))$ when approximating $-1/y$.
This residual is orthogonal to every polynomial of degree less than $r$:
for an odd monomial divide it by $y$ and use orthogonality of $U_r$;
for an even monomial parity makes the integral zero. Therefore
\begin{equation}\tag{FW31}
 H_{2j}^\circ(0)=H_{2j+1}^\circ(0)
                  =M_{2j}^\circ/|U_{2j}^\circ(0)|^2.
\end{equation}
The equality at $2j+1$ also follows because $C_{2j}^\circ(0)=0$.

For completeness the return to the actual complex $z$ has the finite
estimate
\begin{equation}\tag{FW32}
 \left|\log\frac{H_r^\circ(z)}{H_r^\circ(0)}\right|
 \le 2(r+1)[-\log(1-\varrho)]+64(1+R/d)\varrho,
 \quad0\le r\le q+1.
\end{equation}
For $r=0$, use $H_0(z)=\operatorname{Im}[-zM_0 A_0]/d$.
For even $r\ge2$, write (FW6) as $|C_{r-1}|^2$ times the imaginary
part of the first line of (FW14). The factor $C_{r-1}$ is odd in its
polynomial index, so (FW13) cancels its single factor $z$; its ratio
to the value at zero is $A_{r-1}$ times the reciprocal of the individual
squared-root product. Each root factor costs at most
$-\log(1-\varrho)$ in absolute log modulus. The imaginary-part ratio
costs at most $16(R/d)\varrho$ on the guard. For odd $r$, the even
$C_{r-1}$ contributes $|z|^2$ and the second line of (FW14) contributes
$\operatorname{Im}(1/z)=d/|z|^2$; these cancel exactly. The same root
and $A$ estimates apply. The allowance in (FW32) bounds all these
terms, including $r=1$. This argument never treats a polynomial
interpolation residual at nonzero $z$ as an orthogonal projection.

The monic minima and (FW9) also give
$(1-\eta)\nu_j\le\nu_j^\circ\le\nu_j$. Combining (FW10)--(FW11),
(FW26), (FW31)--(FW32) shows
\begin{equation}\tag{FW33}
 \log\frac{E_{q-1}}{E_{2q-1}}
 =\log\nu_0-\log\nu_q+2\log|U_q^\circ(0)|
                     +O_h(1+k^2/q).
\end{equation}
This endpoint computation needs no uniform moving-degree monic estimate.
For each of the two specified monic minima, (FW7) bounds
$\log(\nu_j^m/\nu_j^\Gamma)$ between $\log\ell$ and $\log u$;
their difference is therefore within $\Lambda$. In Gamma, AG19 with
$(n,s)=(q/2,0),(q/2+1,1)$ and AG22 give directly
\begin{equation}\tag{FW34}
 \begin{split}
 \log\nu_q^\Gamma&=4q\log q+(F-L_{\rm eq}+3)q+O_h(\log q),\\
 \log\nu_0^\Gamma&=2q\log q+(4\log2-2\log\pi-2)q+O_h(\log q).
 \end{split}
\end{equation}
Indeed the even relation block acquires just its next monic vector;
its determinant ratio is $H_{q/2+1}^q/H_{q/2}^q$. Subtracting AG19
at these two literal block sizes produces the first line, with its
$qL_{\rm eq}$ term and all powers of $q$ retained.
Thus (FW33) equals
$[2L_{\rm eq}-F-5+4\log2-2\log\pi]q+O_h(k+\log q)$.
The original ECL/LRS16 convention states
$C_\partial=4\log(4/\pi)-10-2F+4L_{\rm eq}$, so this is
$C_\partial q/2+O_h(k+\log q)$. Finally $r=0,q$ are even,
and (FW11),(FW16) bound each corresponding $T_r$ by $O_h(k^2/q^2)$.
Returning through (FW4) proves
\begin{theorem}[Both original endpoint windows]
\begin{equation}\tag{FW35}
 \log\frac{E_{q-1}}{E_{2q-1}}
 =\frac{C_\partial q}{2}+O_h(k+\log q),\qquad
 \log\frac{E_q}{E_{2q}}
 =\frac{C_\partial q}{2}+O_h(k+\log q).
\end{equation}
In particular
\begin{equation}\tag{FW36}
 \boxed{\log\frac{E_{q-1}E_q}{E_{2q-1}E_{2q}}
                    =C_\partial q+O_h(k+\log q).}
\end{equation}
\end{theorem}
All occurrences of $|b_\lambda|^2$ cancel between the same original
class at different cutoffs; they were retained in (FW4) before that
cancellation. This proof evaluates the actual original-source class.

\section{Exact receiving map into the original observation quotient}
Retain the actual fixed observation $\Lambda:E\to B=\operatorname{im}\Lambda$,
including its period map, and an injective frame $I$ for $K=\ker\Lambda$.
Define
\begin{equation}\tag{FW37}
 P_N=I(I^*G_NI)^{-1}I^*G_N,\quad
 Q_N=(\Lambda G_N^{-1}\Lambda^*)^{-1},\quad
 \theta_N=v^*G_NP_Nv/E_N.
\end{equation}
The exact minimum lift $G_N^{-1}\Lambda^*Q_N:B\to E$ is a right
inverse of $\Lambda$, is orthogonal to $K$, and proves
\begin{equation}\tag{FW38}
 G_N(I-P_N)=\Lambda^*Q_N\Lambda,\qquad
 e_{B,N}:=(\Lambda v)^*Q_N\Lambda v=E_N(1-\theta_N).
\end{equation}
This supplies the typed metric map between the just evaluated class norm
and its observed norm. Its complete proof is the same quadratic minimum
as ORE2--3 \cite{ore}; the general minimum framework is credited there
to Boyd and Vandenberghe \cite[Appendix C.4.1]{Boyd}.

If $\Lambda v\ne0$, strict positivity of $Q_N$ gives $1-\theta_N>0$
at every finite cutoff. Equations (FW35)--(FW38) then yield the new
receiving formula
\begin{equation}\tag{FW39}
 \boxed{\log\frac{e_{B,q-1}e_{B,q}}{e_{B,2q-1}e_{B,2q}}
 =C_\partial q+
 \log\frac{(1-\theta_{q-1})(1-\theta_q)}
               {(1-\theta_{2q-1})(1-\theta_{2q})}
 +O_h(k+\log q).}
\end{equation}
The unresolved part in this specific receiver is now four actual endpoint
projection angles. It has not been replaced by ranks or by a scalar total.
If $\Lambda v=0$, the observed energy is identically zero and (FW38),
rather than an undefined logarithm, is the receiving statement.

There is also a positive exact calculation at each original increment.
With $b=b_{N+1}$ and $\beta=b^*G_Nb$, put
$\alpha_N=b^*G_NP_Nb/\beta$,
$d_{K,N}=\det(I^*G_{N+1}I)/\det(I^*G_NI)$, and
$V_N=b^*G_NP_Nv/(b^*G_Nv)$. The original source increment and its
restriction, checked by multiplying the proposed inverse, give
\begin{equation}\tag{FW40}
 G_{N+1}^{-1}=G_N^{-1}+bb^*/\omega_{N+1},\quad
 d_{K,N}=1-(1-\Delta_r)\alpha_N,\quad
 \frac{e_{B,N+1}}{e_{B,N}}
 =1-\frac{T_r|1-V_N|^2}{d_{K,N}(1-\theta_N)}.
\end{equation}
The quotient update denominator is exactly
$\omega_{N+1}+\beta(1-\alpha_N)=(\omega_{N+1}+\beta)d_{K,N}$;
pairing the inverse update with $\Lambda v$ proves (FW40).
The matrix inverse method is Hager's rank-one update \cite{Hager},
used with the original vector and denominator. Combining with
$E_{N+1}/E_N=1-T_r$ gives the undivided identity
\[
 (1-T_r)(1-\theta_{N+1})
 =1-\theta_N-T_r|1-V_N|^2/d_{K,N}.
\]
Thus the evaluated contraction (FW36) and all actual observation responses
belong to this same sequence of attained minima. For any positive sequence
$a_N$, cancellation proves
$\log(a_{q-1}a_q/(a_{2q-1}a_{2q}))=-\sum w_r\log(a_{N+1}/a_N)$.
Applying it to (FW40) recovers (FW39) with the exact endpoint weights.
It retains the finite common scalar phase correction in the three signed
currents; their sum still uses $P_K+P_B+P_\times=1$ in the original
response-product notation. No individual signed current is assigned here.

\section{Degree-resolved contraction of the original class}
Fix $s\in(0,1]$, put $r=\lfloor sq\rfloor$, $n=\lfloor r/2\rfloor$,
and retain the equilibrium for
$V_s(x)=(\pi/s)\sqrt x-(2/s)\log x$.
Write $L_s=\mathcal L(2/s,\pi/s)$, $F_s=-\min\mathcal E_{2/s,\pi/s}$
and $\lambda_s=\lambda(2/s,\pi/s)$, with the original Robin convention
$V_s-2\int\log|x-y|\dd\rho_s(y)=\lambda_s$ on its support.
The comparison proof (FW20)--(FW23) applies to this $n$ without change:
its common spaces still have codimension one, all output degrees remain
at most $2q+2$, and its rational comparison retains the same $B_m,B_\Gamma$.
Choose the fixed gap below the support of this equilibrium as well.
For fixed $s$, this is a change of an auxiliary estimate only; (FW10)
returns its explicitly bounded contribution to the same original minimum.

For the Gamma calculation let $b_s=sq-2n\in[0,2)$. Its exact additional
analytic factor, on top of $-3\log x/4-\tau_k/x$, is
\[
 \frac{b_s}{s}\log x-\frac{\pi b_s}{2s}\sqrt x.
\]
It is retained in both Gram determinants. The analytic-statistic estimate
used in (FW25) is uniform for this bounded pair $(b_s,\tau_k)$: the
logarithmic determinant minus its linear statistic is convex in the pair;
bound it at the vertices and centre of a fixed larger rectangle and use
the same convex-combination argument as for (FW25). Thus
\begin{equation}\tag{FW41}
 \log|U_{2n}^\circ(0)|=2n\log q+nL_s+O_{h,s}(k+\log q).
\end{equation}
Both complementary-region bounds used in AG have the same proof here,
with $D/q\le1+o(1)$ and a fixed positive $s$. Their constants may depend
on $s$; no uniformity as $s\downarrow0$ is used in (FW41).

The original Gamma monic norm of degree $2n$ is the ratio of its even
Gram blocks of sizes $n+1,n$. To use a fixed potential in both, the
larger block's analytic factor is the smaller one's factor plus $V_s$.
The leading log determinant difference is therefore
$2nF_s+n\int V_s\dd\rho_s=-n\lambda_s$; the equality follows by
integrating the Robin equation. The exact row, column and mass powers
of $q$ contribute $(2q+4n+1/2)\log q$.
Comparison of the actual monic minima by (FW7) consequently gives
\begin{equation}\tag{FW42}
 \log\nu_{2n}=(2q+4n)\log q-n\lambda_s+O_{h,s}(k+\log q).
\end{equation}
The term $(1/2)\log q$ has only been absorbed into the stated logarithmic
error, not multiplied by a Gram rank. Combine (FW41)--(FW42) with
(FW31)--(FW32), and use
$\log\nu_0=2q\log q+(2\log(4/\pi)-2)q+O_h(k+\log q)$ from (FW34).
If $r$ is odd, its extra even-index update changes the logarithm by
$O_h(k^2/q^2)$, by (FW16). We obtain, at every fixed multiplicity,
\begin{equation}\tag{FW43}
 \boxed{\begin{gathered}
 -\log\frac{E_{q-1+\lfloor sq\rfloor}}{E_{q-1}}
 =q\mathcal J(s)+O_{h,s}(k+\log q),\\
 \mathcal J(s)=s(L_s+\lambda_s/2)+2a_0-2,\qquad a_0=\log(4/\pi).
 \end{gathered}}
\end{equation}

This is the original degree profile, with an exact expression for its
derivative. In the conventions of LRS10--15 \cite{lrs}, let
\[
 \psi(t)=(1+t)(1-\log(1+t))-\tfrac t4\lambda_t,\qquad
 f(t)=2\psi(t)-2(t+1)\psi'(t).
\]
For the original elliptic modulus $\kappa_t$, write
$r_t=\sqrt{1-\kappa_t^2}$. The endpoint equation and its consequences are
\begin{equation}\tag{FW44}
 \frac{E(\kappa_t)}{r_tK(\kappa_t)}=1+t,\quad
 \psi(t)=\log\frac{1+r_t}{E(\kappa_t)}
                 +t\log\frac{\kappa_t}{E(\kappa_t)},\quad
 f(t)=\log\frac{1+r_t}{1-r_t}>0.
\end{equation}
Here $K,E$ are complete elliptic integrals; the modulus is unrelated to
the observation kernel. The derivative identity
$\psi'=\log(\kappa_t/E(\kappa_t))$ follows from the original endpoint
variation, as proved in LRS8--12. Its convergent small-modulus expansion
gives $\kappa_t=2t^{1/4}(1+O(\sqrt t))$, hence
$f(t)=-\tfrac12\log t+O(\sqrt t)$, an integrable function at zero.
No asymptotic remainder has been differentiated.

For an explicit check of the integrated identity, LRS15 gives the
continuous primitive
$\mathfrak F(t)=t^2F_t/4+(1+t)^2(3/2-\log(1+t))/2-3/4$
of $\psi$, with value zero at zero. Integration by parts gives
\[
 \int_0^s f(t)\dd t=4\mathfrak F(s)-2(s+1)\psi(s)+2a_0.
\]
Substitution and the proved equilibrium equality
$F_s=-\lambda_s/2-(2/s+1)+L_s/s$ reduce this expression exactly to
$s(L_s+\lambda_s/2)+2a_0-2$. Thus $\mathcal J(s)=\int_0^s f$,
including its continuous zero endpoint.

Put $g_r=H_{r+1}/H_r=1-T_r$. The finite positive measures
\begin{equation}\tag{FW45}
 \mu_k=\frac1q\sum_{\substack{0\le r<q\\r\ {
m odd}}}
              (-\log g_r)\delta_{r/q}
 \quad\hbox{converge weakly on $[0,1]$ to }f(s)\dd s.
\end{equation}
Indeed (FW16) bounds the sum of all even-index losses by
$O_h(k^2/q)=o(q)$. The full prefix sum is exactly
$-\log(H_{\lfloor sq\rfloor}/H_0)$, whose quotient by $q$ tends to
$\mathcal J(s)$ by (FW43). For any fixed $s>0$, sandwiching its boundary
atom between prefixes at $s-\epsilon$ and $s+\epsilon$ proves convergence
at $s$ by continuity of $\mathcal J$. At zero, total prefix mass is at
most the mass below any positive $\epsilon$, whose limit tends to zero.
At one, (FW35) gives the full mass. These statements prove weak convergence
against continuous functions by step-function approximation. They do not
claim a pointwise limit of any individual odd removal fraction.

\section{Original conductor bound at every cutoff and its determinant return}
This section retains the original simple-quartet conductor domain:
$m_0=1$, $k\ge9$, its fixed actual period $\varpi$, and the finite root
guards in FEX/LRS. Its full original observation has
$K=\ker\Lambda_k\subset\{[P]:\chi_{k-8}\mid\mathcal T_AP\}$, with
\begin{equation}\tag{FW46}
 \begin{gathered}
 q'=(k-7)^2,\quad \Delta_A=q-q'=16k-48,\quad c'=c-4,\\
 \mathcal T_AP(S')=\sum_{a,b=0}^{8}a_{ab}P(S'+\beta_{8,a,b}),
 \quad \beta_{8,a,b}=4+(2a-8)\delta+i(2b-8)\gamma,\\
 \mu_j=\sum a_{ab}\beta_{8,a,b}^j,\quad
 v=\min\{j:\mu_j\ne0\}\le80,\quad
 \mathcal C_A=(v!/\mu_v)\mathcal T_A.
 \end{gathered}
\end{equation}
All its constants and its scalar are the original ones, as reproduced
in the full conductor proof \cite{oid}. At a lower root, every shifted
point is an upper root. Consequently $\mathcal T_A(\chi H)$ vanishes
at every lower root and is divisible by the complete simple lower
polynomial. Thus the conductor map is defined on the original quotient
and factors through the actual observation; this proves its use on every
minimum representative, rather than only on a chosen coefficient section.
The leading coefficient of $\mathcal C_AP$ in degree $N-v$ is
$(N)_v$ times the leading coefficient of $P$, by the finite Taylor formula.

The forward Gamma bound retains the centre displacement exactly:
\[
 \norm{\mathcal C_AP}_{\sigma,c-4}\le M_N\norm P_{\sigma,c},\quad
 M_N=\max\left\{1,\frac{v!}{|\mu_v|}\sum|a_{ab}|
             e^{|\beta_{8,a,b}-4|(\mathsf H_N+2)}\right\},\quad
 \mathsf H_N=\sum_{j=1}^N\frac1j.
\]
The full derivation in \cite{oid} expands differentiation in the
Meixner--Pollaczek basis and bounds its weighted row and column sums
by $\mathsf H_N$ and $\mathsf H_N+4$, respectively. Exponentiating the
finite nilpotent derivative gives exactly this translation bound at every
finite degree $N$, so $\log M_N=O_{h,\varpi}(\log q)$ on the window.

Let $\nu^-_{j,1}$ be the full monic lower-relation minimum in the
order-one Gamma source. For an arbitrary $x\in K$, write its actual
minimum representative as $T_Nx=P$ and let its leading coefficient be
$a_N(P)$. The divisibility just proved and (FW7) imply
\[
 \norm x_{G_N}^2\ge\ell M_N^{-2}(N)_v^2
                 \nu^-_{N-v-q',1}|a_N(P)|^2.
\]
This is also true if the leading coefficient is zero. The exact boundary
pairing is $b_N^*G_Nx=\omega_Na_N(P)$.
Take its squared dual norm on the entire $K$, use
$\omega_N\le u\gamma_N$, and retain
$F_{11}=\omega_N\eta_N$, with $\eta_{q-1}=1$ and
$\eta_N=1-\Delta_{N-q}$ for $N\ge q$. This proves, throughout the window,
\begin{equation}\tag{FW47}
 \pi_{1,K,N}:=\frac{\norm{P_Nb_N}_{G_N}^2}{F_{11}}
 \le\min\left\{1,\frac{u}{\ell}
 \frac{M_N^2\gamma_N}{\eta_N(N)_v^2\nu^-_{N-v-q',1}}\right\},
 \quad q-1\le N\le2q.
\end{equation}
Its lower-relation degree is nonnegative on the original eventual guards.
The inherited LRS1 proves the upper degree bound $\le2q'$ there.
All finite lower endpoints LRS4 may be inserted for the displayed monic
norm before taking asymptotics; hence (FW47) has an explicit finite
receiver containing no observation determinant.

For the rate use LRS3--11, in the order-one reference for both complete
simple root lists. It gives, uniformly for $0\le j\le2Q$,
\[
 \begin{gathered}
 \log(\nu_{j,1}^{(b)}/\gamma_{Q+j})=2Q\psi(j/Q)+O_h(\log q),\\
 Q=(b+1)^2,\qquad b=k,k-8.
 \end{gathered}
\]
The original arithmetic monic norms and relation norms differ from their
Gamma values by log factors in $[\log\ell,\log u]$; thus their ratio
costs at most $\Lambda=O_h(k+\log q)$. The exact boundary-radius formula
preceding (FW19), (FW16), and the current identity used in (FW28) now give
\begin{equation}\tag{FW48}
 -\log p_{1,q-1+r}=2q\psi(r/q)+O_h(k+\log q),
 \quad 0\le r\le q+1.
\end{equation}
For clarity, $\log\kappa_N^2$ equals
$\log(\nu_r/\omega_{q+r})+\log(\omega_{N+1}/\omega_N)$ plus the
two actual $\log(1-\Delta)$ factors. Comparison with
$\gamma_{n+1}/\gamma_n=(n+1)(n+1/2)$ bounds that recurrence logarithm
by $O_h(k+\log q)$; (FW19) bounds the remaining factors. Also (FW11),
(FW16) and $p_2\le1$ give $|\log p_2|=O_h(\log q)$.
These observations prove (FW48) including $r=0$.

Set $F_*(n,r)=2n\psi(r/n)$ to avoid confusion with the equilibrium
constant $F$. Its exact derivative difference is
$\partial_nF_*-\partial_rF_*=f(r/n)$, by (FW44). Therefore
\[
 F_*(q,r)-F_*(q',r+\Delta_A)
 =\int_0^{\Delta_A}f\left(\frac{r+s}{q-s}\right)\dd s.
\]
The integral at $r=0$ is convergent. A shift of the second argument
by the fixed $v+1$ costs $O_h(\log q)$ using the proven logarithmic
endpoint bound for $\psi'$. Since $f(t)\le-\frac12\log t+C$
for $t\le1$ and is bounded on the remaining compact interval, its
integral is at most
$\frac{\Delta_A}{2}\log^+(q/(r+\Delta_A))+C\Delta_A$.
The bound follows by integrating $\log(q/(r+s))$ exactly; its excess
over $\Delta_A\log(q/(r+\Delta_A))$ is
$\Delta_A-r\log(1+\Delta_A/r)\le\Delta_A$.
Finally the monic factorial factor satisfies the exact identity
\[
 \frac{\gamma_N}{\gamma_{N-v}(N)_v^2}
       =\prod_{j=0}^{v-1}\frac{N-j-1/2}{N-j}\le1.
\]
Using it in (FW47) and then dividing by (FW48) gives
\begin{equation}\tag{FW49}
 |U_N|\le\sqrt{\theta_N}
 \exp\left\{\frac{\Delta_A}{4}\log^+\frac{q}{r+\Delta_A}
                    +C_{h,\varpi}k\right\},\qquad r=N+1-q.
\end{equation}
Indeed $|U_N|^2\le\theta_N\pi_{1,K,N}/p_{1,N}$ is precisely
Cauchy--Schwarz on the two actual $K$ projections. Similarly
$|V_N|^2\le\theta_N/p_{2,N}$ and (FW16) yield
\begin{equation}\tag{FW50}
 |V_N|\le C_h\sqrt{\theta_N}
 \begin{cases}q/k,&r\text{ even},\\1,&r\text{ odd}.
 \end{cases}
\end{equation}
For odd $r$, its lower bound follows explicitly from
$b_r\ge Y^2/\widehat C^2$; no gap below one is asserted.

There is a further direct determinant receiver. The original rank-one
update (FW40), now applied to $b=b_{N+1}$ in $K$, gives
$\pi_{1,K,N+1}=\Delta_r\alpha_N/d_{K,N}$.
Solving this equality with the displayed value of $d_{K,N}$ yields
\begin{equation}\tag{FW51}
 d_{K,N}=\frac{\Delta_r}{\Delta_r+(1-\Delta_r)\pi_{1,K,N+1}},\qquad
 0\le-\log d_{K,N}
 \le\frac{\Delta_A}{2}\log^+\frac{q}{r+\Delta_A}+C_{h,\varpi}k.
\end{equation}
For the inequality repeat (FW47) at $N+1$ and compare its monic norm
with $-\log\Delta_r=2q\psi(r/q)+O_h(k+\log q)$.
The lower index is now $r+\Delta_A-v$, so the identical integral estimate
applies with the fixed shift $v$. The logarithm $\log(1+e^x)$ costs at
most $\log2+\max(x,0)$; all constants are absorbed in the specified
$C_{h,\varpi}k$. This is an inequality for the same original restriction
determinant, rather than an allocation from its rank.

In the complete weighted window, comparison with
$\int_0^q\log(q/x)\dd x=q$ gives
\begin{equation}\tag{FW52}
 \boxed{0\le
 \log\frac{\det(I^*G_{q-1}I)\det(I^*G_qI)}
               {\det(I^*G_{2q-1}I)\det(I^*G_{2q}I)}
 =\sum_{r=0}^q w_r[-\log d_{K,q-1+r}]
 \le C_{h,\varpi}'kq.}
\end{equation}
This carries the established $kq$ control into the literal observation
kernel through its actual conductor. It does not equate this object with
another native flag. The exact morphism to the observation quotient is
$\det Q_{N+1}/\det Q_N=\Delta_r/d_{K,N}$ from (FW40); its full-window
loss is the already evaluated original total minus the exact restriction
loss in (FW52). The latter bound does not specify its order-$kq$ coefficient.

\section{Precise integration scope}
Equations (FW27)--(FW30) replace the low-cutoff-only source-product status
in the preceding edition. Equations (FW35)--(FW36) evaluate its previously
unevaluated full terminal-class norm return. Equations (FW38)--(FW40)
push those results into ORE's original observation, including the full
conductor cross blocks retained in COI11--13 and ORE19. The exact finite
term $D_k$ in (FW18) improves the incoming proof's dependency: an
exponential estimate for $\Delta_r$ is unnecessary for those conclusions.
Equations (FW41)--(FW45) prove the further fixed-degree-fraction profile
and weak loss measure. Equations (FW46)--(FW52) prove the complete-window
conductor response bound on its original simple-quartet domain and carry
it into the original restriction determinant. The uniform moving-degree
first-energy estimate (FW48) is certified here on that same simple domain;
the main product and class-profile theorems retain every fixed multiplicity.
The seven
native determinant allocations and projected signs retain their actual
period-dependent projections. No claim to an RH proof is made.

\begin{thebibliography}{99}
\bibitem{incoming} User-supplied mathematical continuation,
\emph{The calculation now extends from the two low cutoffs to the complete
original window}, 19 September 2026, equations (1)--(46).
The two deliveries are byte-identical; the original Markdown and SHA256
are preserved in this edition's input receipt.
\bibitem{DLMF} T.~H.~Koornwinder, R.~Wong, R.~Koekoek,
R.~F.~Swarttouw and W.~P.~Reinhardt, Chapter 18,
\emph{NIST Digital Library of Mathematical Functions},
\S18.2(iv), (vi), (vii), \url{https://dlmf.nist.gov/18.2}.
Used for the classical recurrence, interlacing and quadratic transformation.
\bibitem{BG} Ga\"etan Borot and Alice Guionnet,
``Asymptotic expansion of $\beta$ matrix models in the one-cut regime'',
\emph{Communications in Mathematical Physics} 317 (2013), 447--483,
\url{https://arxiv.org/abs/1107.1167}, Propositions 1.2 and 5.2.
Used for the fixed-positive-interval analytic-statistic expansion in (FW25).
\bibitem{Boyd} Stephen Boyd and Lieven Vandenberghe,
\emph{Convex Optimization}, Cambridge University Press, 2004,
Appendix C.4.1, \url{https://web.stanford.edu/~boyd/cvxbook/bv_cvxbook.pdf}.
\bibitem{Hager} William W.~Hager, ``Updating the Inverse of a Matrix'',
\emph{SIAM Review} 31 (1989), 221--239,
\url{https://doi.org/10.1137/1031049}, equation (2).
\bibitem{cai} Split-Zero programme, \emph{Complete canonical action
prerequisite}, CAI16--28, full cited provider retained in this edition.
\bibitem{rbe} Split-Zero programme, \emph{The retained eigenclass's two
original boundary energies}, 19 September 2026, RBE1--29.
\bibitem{ag} Split-Zero programme, \emph{The absolute Gamma baseline and
the complete action at logarithmic precision}, 19 September 2026,
AG1--29. Its complete proof and source providers are retained.
\bibitem{ore} Split-Zero programme, \emph{Adjacent original observation
responses and their exact metric evolution}, 19 September 2026, ORE1--23.
\bibitem{lrs} Split-Zero programme, \emph{Literal covariance, logarithmic
endpoint term, and the complete action}, 18 September 2026, LRS1--29,
\texttt{LRC\_SCALAR\_RECONSTRUCTION.tex}. Complete root and finite norm
providers LRC11--15, EQ12--36 and EIQ/ECL are retained in the cumulative source.
\bibitem{oid} Split-Zero programme, \emph{Full source derivation and the
first conductor allocation}, 19 September 2026,
\texttt{ORE\_INDEPENDENT\_DERIVATION.tex}, section ``The actual conductor
evaluates the first logarithmic allocation''. This source reproduces the
original FEX conductor, its finite forward bound and its exact quotient map.
\end{thebibliography}
\end{document}
